\section*{Orientaciones bibliográficas}
\addcontentsline{toc}{section}{Orientaciones bibliográficas}
Para submódulos, cocientes y construcciones libres, véase \citet[cap.~III, 5.8--18 y apartados 6.1--3]{aluffi2009}. La acción del cociente corresponde a III.5.11 y la inyectividad de los generadores es análoga a la comprobación II.5.3.

Para noetherianidad y álgebras finitas, véase \citet[cap.~III, 6.4--9]{aluffi2009}; para Hilbert e irreducibilidad, el cap.~V, 1.1--6. Los ejercicios corresponden a III.6.14, V.1.6 y V.1.12.

Para bases, rango y torsión, véase \citet[cap.~VI, 1.1/4--7/9--11/13--14 y 4.1--3]{aluffi2009}. El lema de separación es VI.5.2 y los ejercicios corresponden a VI.1.8 y VI.4.1.

Para la clasificación de módulos sobre un DIP, véase \citet[cap.~VI, 5.1--7]{aluffi2009}; los módulos cíclicos y anuladores se definen en VI.4.4/6. Las comprobaciones de rango y anuladores corresponden a los ejercicios VI.5.6 y VI.4.4.

Para módulos, álgebras y homomorfismos, véase \citet[cap.~III, 5.1--7]{aluffi2009}. Los ejemplos de factorización posteriores a Hilbert se relacionan con V.1.7--12 y V.4.10--18; el de infinitas variables corresponde al ejercicio V.1.15.
